\documentclass[11pt,a4paper]{article}
\usepackage[margin=2.3cm]{geometry}
\usepackage{graphicx}
\usepackage{amsmath,amssymb}
\usepackage{booktabs}
\usepackage{tabularx}
\usepackage{array}
\usepackage{multirow}
\usepackage{float}
\usepackage{hyperref}
\usepackage{bookmark}
\usepackage{xcolor}
\usepackage{fontspec}
\usepackage{ucharclasses}

\setmainfont[ItalicFont={Noto Sans CJK KR}]{Noto Sans CJK KR}
\newfontfamily\hangulfont[ItalicFont={Noto Sans CJK KR}]{Noto Sans CJK KR}
\setTransitionsForCJK{\hangulfont}{}{}

\tolerance=1000
\emergencystretch=3em
\hfuzz=2pt

\definecolor{tableheader}{HTML}{E3F2FD}
\definecolor{lightgray}{HTML}{F5F5F5}
\definecolor{goodblue}{HTML}{1565C0}
\definecolor{warnred}{HTML}{C62828}
\definecolor{okgreen}{HTML}{2E7D32}

\title{\textbf{Exp 4-2 추가 실험 별도 보고서}\\[0.4em]
\large WikiText-2 Sliding-Window PPL Benchmark 결과 정리}
\author{Boltzmann Attention / FOKVQ}
\date{2026년 4월 1일}

\begin{document}
\maketitle

\section{개요}

본 보고서는 \texttt{gpt2-medium} 기반 \texttt{WikiText-2} sliding-window perplexity 벤치마크에서 수행한 \textbf{Exp 4-2 추가 실험}의 최종 정리본이다. 목적은 세 가지였다. 첫째, \texttt{FOKVQ}가 공정한 평균 비트 예산 하에서 \texttt{uniform}, \texttt{KIVI-style}, \texttt{turboquant-style}보다 우수한지 확인한다. 둘째, 저비트에서 \texttt{PCA basis}가 \texttt{random} 또는 \texttt{identity} basis보다 실제로 유리한지 확인한다. 셋째, 이러한 경향이 더 긴 문맥에서도 유지되는지 점검한다.

핵심 결론은 다음과 같다. \texttt{FOKVQ}는 현재 구현과 프로토콜에서는 \texttt{KIVI-style} 및 \texttt{turboquant-style}보다 우수하지 않다. \texttt{4-bit}에서는 거의 모든 방법이 근접하여 \texttt{FOKVQ}의 이점이 드러나지 않았고, \texttt{2-bit}에서는 \texttt{PCA basis}가 \texttt{random}/\texttt{identity}보다 낫지만 strong baseline을 넘지는 못했다. 따라서 현재 결과는 우월성(superiority) 증거라기보다 \textbf{방법론적이고 부정적이지만 유의미한 결과}로 해석하는 편이 정확하다.

\section{실험 프로토콜}

\begin{table}[H]
\centering
\caption{공통 실험 설정}
\label{tab:setup}
\small
\begin{tabularx}{\textwidth}{@{}lX@{}}
\toprule
\textbf{항목} & \textbf{설정} \\
\midrule
모델 & \texttt{openai-community/gpt2-medium} \\
데이터셋 & \texttt{WikiText-2} \\
평가 방식 & sliding-window causal LM perplexity \\
최대 평가 토큰 수 & \texttt{4096} \\
Calibration samples & \texttt{8} \\
Attention implementation & \texttt{eager} \\
시드 & \texttt{42} \\
비교 방법 & \texttt{uniform}, \texttt{KIVI-style}, \texttt{turboquant-style}, \texttt{FOKVQ(PCA)} \\
주의사항 1 & \texttt{turboquant-style}은 random orthogonal rotation + Lloyd-Max scalar quantization 기반 내부 구현이며, 공식 TurboQuant 완전 재현으로 해석하면 안 된다. \\
주의사항 2 & \texttt{FOKVQ} 결과는 공정한 평균 비트 예산을 강제한 이후의 실험만 반영한다. \\
\bottomrule
\end{tabularx}
\end{table}

\section{핵심 결과}

\subsection{4-bit 메인 비교}

\begin{table}[H]
\centering
\caption{4-bit main comparison (\texttt{context\_len=256}, \texttt{stride=128})}
\label{tab:4bit-main}
\small
\begin{tabular}{lr}
\toprule
\textbf{Method} & \textbf{PPL} \\
\midrule
FP16 & 26.8638 \\
KIVI-style & 26.8770 \\
uniform & 26.8983 \\
FOKVQ & 26.9954 \\
turboquant-style & 27.0020 \\
\bottomrule
\end{tabular}
\end{table}

이 표의 해석은 단순하다. \texttt{4-bit}에서는 모든 방법이 매우 가깝고, \texttt{FOKVQ}는 \texttt{KIVI-style}보다 낫지 않다. 따라서 \texttt{4-bit}는 현재 구현의 주력 주장 영역이 아니다.

\subsection{4-bit 축 ablation}

\begin{table}[H]
\centering
\caption{4-bit axis ablation (\texttt{context\_len=256}, \texttt{stride=128})}
\label{tab:4bit-axis}
\small
\begin{tabular}{lr}
\toprule
\textbf{Method} & \textbf{PPL} \\
\midrule
FP16 & 26.8638 \\
uniform & 26.8983 \\
random basis & 26.9608 \\
FOKVQ(PCA) & 26.9954 \\
identity basis & 27.0862 \\
\bottomrule
\end{tabular}
\end{table}

\texttt{identity}는 가장 나쁘지만, \texttt{random}이 \texttt{PCA}보다 약간 더 좋다. 즉 \texttt{4-bit}에서는 \texttt{PCA basis}의 구조적 우위가 관찰되지 않는다.

\subsection{2-bit 공정 비교: \texttt{topk=0.25} vs \texttt{topk=0.5}}

\begin{table}[H]
\centering
\caption{2-bit fair main comparison}
\label{tab:2bit-main}
\small
\begin{tabular}{lrr}
\toprule
\textbf{Method} & \textbf{topk=0.25} & \textbf{topk=0.5} \\
\midrule
FP16 & 26.8638 & 26.8638 \\
KIVI-style & 27.7507 & 27.7507 \\
turboquant-style & 27.7524 & 27.7524 \\
uniform & 32.4926 & 32.4926 \\
FOKVQ & 48.2813 & 34.5466 \\
\bottomrule
\end{tabular}
\end{table}

\texttt{topk=0.5}는 \texttt{0.25}보다 현저히 낫다. 그러나 \texttt{topk=0.5}에서도 \texttt{FOKVQ}는 여전히 \texttt{uniform}, \texttt{KIVI-style}, \texttt{turboquant-style}보다 좋지 않다.

\begin{figure}[H]
\centering
\includegraphics[width=0.92\textwidth]{figures/exp4_2_2bit_topk050_contexts.png}
\caption{2-bit 공정 비교에서 \texttt{topk=0.5} 설정의 문맥 길이별 PPL 변화. \texttt{FOKVQ}는 긴 문맥으로 갈수록 개선되지만, 모든 구간에서 \texttt{turboquant-style} 및 \texttt{KIVI-style}보다 열세이다.}
\label{fig:2bit-context}
\end{figure}

Figure~\ref{fig:2bit-context}에서 \texttt{FOKVQ(topk=0.5)}는 문맥이 길어질수록 절대 PPL이 낮아진다. 그러나 같은 구간에서 \texttt{turboquant-style}과 \texttt{KIVI-style}도 더 좋아지므로 순위는 바뀌지 않는다. 즉 긴 문맥이 \texttt{FOKVQ}에만 유리하게 작동하지는 않았다.

\subsection{2-bit 장문맥 비교}

\begin{table}[H]
\centering
\caption{2-bit fair comparison across longer contexts}
\label{tab:2bit-long}
\small
\begin{tabular}{lrr}
\toprule
\textbf{Method} & \textbf{512 / 256} & \textbf{1024 / 512} \\
\midrule
FP16 & 22.5366 & 18.3817 \\
turboquant-style & 23.3546 & 19.3868 \\
KIVI-style & 23.9904 & 20.4618 \\
uniform & 27.3654 & 22.4624 \\
FOKVQ (\texttt{topk=0.25}) & 45.9687 & 42.2597 \\
FOKVQ (\texttt{topk=0.5}) & 28.8335 & 24.0110 \\
\bottomrule
\end{tabular}
\end{table}

\texttt{topk=0.25}는 장문맥에서도 회복되지 않는다. \texttt{topk=0.5}는 훨씬 개선되지만, 여전히 strong baseline보다 낮은 순위에 머문다. 특히 \texttt{1024/512}에서 \texttt{turboquant-style}이 19.3868로 가장 강한 양자화 baseline이다.

\subsection{축 선택 효과: PCA vs random vs identity}

\begin{table}[H]
\centering
\caption{Axis ablation summary}
\label{tab:axis-summary}
\small
\begin{tabular}{lrrr}
\toprule
\textbf{Setting} & \textbf{identity} & \textbf{random} & \textbf{FOKVQ(PCA)} \\
\midrule
2-bit, \texttt{topk=0.25} & 412.4942 & 84.4536 & 48.2813 \\
2-bit, \texttt{topk=0.5} & 149.8334 & 52.9955 & 34.5466 \\
4-bit & 27.0862 & 26.9608 & 26.9954 \\
\bottomrule
\end{tabular}
\end{table}

\begin{figure}[H]
\centering
\includegraphics[width=0.9\textwidth]{figures/exp4_2_axis_ablation.png}
\caption{축 선택 ablation. 저비트에서는 \texttt{PCA basis}가 \texttt{random} 및 \texttt{identity}보다 확실히 낫지만, \texttt{4-bit}에서는 이 차이가 사실상 사라진다.}
\label{fig:axis-ablation}
\end{figure}

Figure~\ref{fig:axis-ablation}과 표~\ref{tab:axis-summary}은 \texttt{PCA basis}의 역할을 가장 명확히 보여준다. \texttt{2-bit}에서는 \texttt{PCA}가 \texttt{random} 및 \texttt{identity}보다 유의미하게 낫다. 따라서 \textbf{축 선택 효과 자체는 실재한다}. 다만 이 효과만으로 strong baseline을 넘지는 못한다.

\section{결과 해석}

\subsection{주장 가능한 것}

\begin{enumerate}
\item \texttt{FOKVQ}에는 저비트에서 의미 있는 \texttt{basis-selection effect}가 있다.
\item \texttt{PCA basis}는 naive \texttt{random} 또는 \texttt{identity} basis보다 낫다.
\item 공정한 평균 비트 예산을 강제하면 이전의 일부 개선은 bit-budget artifact였음이 드러난다.
\end{enumerate}

\subsection{주장하면 안 되는 것}

\begin{enumerate}
\item \texttt{FOKVQ}가 \texttt{KIVI-style}보다 우수하다.
\item \texttt{FOKVQ}가 \texttt{turboquant-style}보다 우수하다.
\item \texttt{4-bit}에서 \texttt{PCA basis}가 분명한 우위를 보인다.
\item 현재 \texttt{turboquant-style} 구현이 공식 TurboQuant 재현이다.
\end{enumerate}

\subsection{연구적으로 가지는 의미}

현재 결과는 ``새로운 최고 성능 방법''의 증거라기보다, \textbf{축 정렬이 왜 일부 경우에만 유효하고 strong baseline 우위로 이어지지 않는지}를 보여주는 분석적 결과에 가깝다. 특히 \texttt{2-bit} 축 ablation은 \texttt{PCA}가 무의미하지 않음을 보여주지만, 동시에 그 효과의 크기가 실전적인 baseline 우위를 만들 정도는 아님을 입증한다.

\section{운영상 발견}

병렬 실험 배치에서 GPU보다 CPU thread oversubscription이 더 큰 병목이었다. 다음 환경 제한을 둔 후 wall-clock runtime이 약 \texttt{9--10분}에서 \texttt{6--12초} 수준으로 줄었다.

\begin{itemize}
\item \texttt{OMP\_NUM\_THREADS=8}
\item \texttt{MKL\_NUM\_THREADS=8}
\item \texttt{OPENBLAS\_NUM\_THREADS=8}
\item \texttt{NUMEXPR\_NUM\_THREADS=8}
\end{itemize}

이 운영상 수정은 프로토콜을 바꾸지 않고 실행 효율만 크게 개선했으므로, 이후 원격 실험의 기본 설정으로 유지하는 것이 적절하다.

\section{최종 결론과 다음 단계}

가장 정직한 결론은 다음과 같다.

\begin{quote}
\texttt{FOKVQ}는 저비트에서 \texttt{PCA} 기반 축 정렬 효과를 보이지만, 공정한 비교 조건에서는 \texttt{KIVI-style} 및 \texttt{turboquant-style}보다 우수하지 않다.
\end{quote}

따라서 현재 결과를 논문화한다면 ``새로운 SOTA 방법''보다는 ``축 기반 KV quantization의 조건부 효과와 한계''를 드러내는 보고서로 정리하는 편이 맞다. 다음 실험 우선순위는 다음 세 가지가 적절하다.

\begin{enumerate}
\item layer-wise 또는 head-wise adaptive basis selection
\item PCA coefficient clipping / normalization / robust scaling
\item memory, latency까지 포함한 trade-off 분석
\end{enumerate}

\section*{원본 결과 파일}

\begingroup
\small
\sloppy
\noindent\path{/scratch/amlt_code/outputs/exp4_2_followup_4bit/main/gpt2-medium-4bit-main_standard_ppl.json}\par
\noindent\path{/scratch/amlt_code/outputs/exp4_2_followup_4bit/axis_ablation/gpt2-medium-4bit-axis-ablation_standard_ppl.json}\par
\noindent\path{/scratch/amlt_code/outputs/exp4_2_followup_2bit_topk025/main/gpt2-medium-2bit-topk025-main_standard_ppl.json}\par
\noindent\path{/scratch/amlt_code/outputs/exp4_2_followup_2bit_topk025/axis_ablation/gpt2-medium-2bit-topk025-axis-ablation_standard_ppl.json}\par
\noindent\path{/scratch/amlt_code/outputs/exp4_2_followup_2bit_topk050/main/gpt2-medium-2bit-topk050-main_standard_ppl.json}\par
\noindent\path{/scratch/amlt_code/outputs/exp4_2_followup_2bit_topk050/axis_ablation/gpt2-medium-2bit-topk050-axis-ablation_standard_ppl.json}\par
\noindent\path{/scratch/amlt_code/outputs/exp4_2_followup_2bit_topk050/longctx512/gpt2-medium-2bit-topk050-longctx512_standard_ppl.json}\par
\noindent\path{/scratch/amlt_code/outputs/exp4_2_followup_2bit_topk050/longctx1024/gpt2-medium-2bit-topk050-longctx1024_standard_ppl.json}\par
\endgroup

\end{document}
